\chuongkhongso{TÓM TẮT ĐỒ ÁN}

Đồ án "Phân tích và thiết kế hệ thống quản lý chuỗi nhà trọ" được thực
hiện nhằm xây dựng một nền tảng quản trị tập trung cho các chủ trọ sở
hữu nhiều khu trọ cùng lúc tại các thành phố lớn ở Việt Nam. Trong bối
cảnh quá trình đô thị hoá diễn ra mạnh mẽ, nhu cầu thuê trọ tăng cao,
nhưng đa số chủ trọ vẫn quản lý bằng sổ sách thủ công hoặc các bảng
Excel riêng lẻ cho từng cơ sở. Hệ quả là dữ liệu phân mảnh, sai sót
trong tính toán, khó tổng hợp doanh thu - công nợ và không có kênh tương
tác minh bạch với khách thuê.

Hệ thống được đề xuất giải quyết các vấn đề trên thông qua một kiến trúc
đa người dùng (multi-tenant) cho phép một chủ trọ (kiêm vai trò quản trị
viên hệ thống và kế toán) quản lý nhiều nhà trọ, mỗi nhà trọ có thể phân
công cho một hoặc nhiều quản lý vận hành riêng. Khách thuê và khách vãng
lai có cổng truy cập riêng để tra cứu hoá đơn, tìm phòng và thanh toán
online (qua cổng VNPay Sandbox tích hợp thật hoặc mã VietQR động; định
hướng MoMo ở v2). Toàn bộ luồng nghiệp vụ từ ký hợp đồng, ghi chỉ số
điện nước, xuất hoá đơn cho đến báo cáo doanh thu đều được số hoá và tự
động hoá tối đa.

Đồ án áp dụng phương pháp phân tích thiết kế hệ thống hướng đối tượng,
sử dụng ngôn ngữ mô hình hoá UML với các biểu đồ: Use Case, Activity,
Class và Sequence để mô hình hoá bài toán một cách trực quan và chi
tiết. Báo cáo xác định 4 tác nhân và hơn 30 ca sử dụng được tổ chức
thành 6 nhóm chức năng chính. Sản phẩm cuối cùng đã được hiện thực hoá
thành một hệ thống đầy đủ chạy trên nền tảng Web Single-Page Application
(ReactJS 18 + Vite 5) với backend Node.js Express + MongoDB Atlas
(NoSQL), tích hợp trợ lý ảo AI Google Gemini và OTP qua Gmail SMTP. Hệ
thống hỗ trợ đầy đủ các luồng nghiệp vụ từ đăng ký, ký hợp đồng, ghi chỉ
số điện nước, sinh hoá đơn, thanh toán đến báo cáo doanh thu và công nợ.

